\section{Learning the local law through noise}
\label{sec:learning}

Now drop the assumption that $\kernel$ is known. We observe $\nseq$ noised sequences
$\vx^{\mu}$, each at one level $t_\mu$, never the clean chain, and maximise the marginal
log-likelihood $\Lik(\params)=\sum_\mu\log\Pt[t_\mu,\params](\vx^\mu)$ of a parameterised kernel
$\kernelth$. Because \cref{sec:inference} computes this likelihood and its gradient exactly,
learning the kernel and inferring the posterior are the same recursion.

\paragraph{The exact gradient from one sweep.}
Because only the transition factors depend on $\params$, Fisher's identity
\citep{cappe2005} gives, on the grid,
\begin{equation}
\begin{aligned}
  \nabla_\params\Lik(\params)&=\bigl\langle\Xis(\params),\,\nabla_\params\log\kernelth\bigr\rangle,\\
  \Xis_{kl}&=\sum_{\mu}\sum_{i\ge2}P(a_i=\gridpt{k},a_{i-1}=\gridpt{l}\mid\vx^\mu),
\end{aligned}
\label{eq:fisher}
\end{equation}
where the expected transition mass $\Xis$ is read off the pairwise marginals that the same
forward--backward pass returns. No differentiation through the recursion is needed, and the identity
is verified against finite differences and against brute-force enumeration of the evidence on short
chains (\codefile{tests/test\_em\_bp.py}, \codefile{tests/test\_fisher\_bruteforce.py}). $\Xis$ is also all
the M-step of EM needs, so EM alternates the two halves of \cref{fig:method}: run BP with the current
kernel to get $\Xis$ (E-step), refit the kernel to $\Xis$ (M-step). The estimator sees one frozen
noised view per sequence and never a clean trajectory; each reported fit selects its iteration
count on a validation bundle disjoint from the test bundle, by the held-out risk against the exact score
or, with no true score anywhere, by the held-out one-view log-evidence (\cref{app:protocol}).

\paragraph{The kernel families.}
The M-step depends on how the kernel is parameterised, and the choice is the main modelling decision.
\begin{itemize}[leftmargin=1.2em,itemsep=0.1em,topsep=0.2em]
\item \emph{Mixture kernel} (25 numbers). For the linear chain,
$\kernelth(a'\mid a)=\sum_{c=1}^{\ncomp}\pi_c\,\normal{a'-\corr a-\nu_c}{s_c^2}$: a linear drift with
driving noise an $\ncomp$-component Gaussian mixture. $\ncomp=1$ is the linear-Gaussian kernel and recovers
the Gaussian score (\cref{prop:lmmse}); $\ncomp=8$ is the paper's estimator. For $\corr$ it is the
belief-weighted Yule--Walker pair $\corr=S_{01}/S_{00}$, $\ivar=(S_{11}-\corr S_{01})/W$ with
$S_{00}=\sum_{kl}\Xis_{kl}\gridpt{l}^2$, $S_{01}=\sum_{kl}\Xis_{kl}\gridpt{k}\gridpt{l}$,
$S_{11}=\sum_{kl}\Xis_{kl}\gridpt{k}^2$ and $W=\sum_{kl}\Xis_{kl}$; the complete-data problem for the
mixture carries a second latent label and we take conditional-maximisation sweeps, a generalised EM
whose ascent is monotone \citep{dempster1977maximum,meng1993maximum}.
\item \emph{Tilted-walk kernels} (7, 25 numbers). For a nonlinear drift,
$\kernelth(a'\mid a)=\normal{a';a}{\bondvar}\,e^{B(a')\cdot\params}/Z_\params(a)$, a Gaussian step with a
learned tilt of the arrival point, $B$ a Legendre or cubic-spline basis. It contains the field's kernel
\cref{eq:tilted} exactly, and the M-step objective is concave in $\params$ at fixed $\bondvar$.
\item \emph{The site energy} (5 numbers). Every tilted walk is reversible with the symmetric pair energy
$W(a,a')=(a'-a)^2/2\bondvar+f(a)+f(a')$, so the directed kernel and the pair energy are one model class
(\cref{app:field}). Parameterising $f$ as a quartic polynomial and $\bondvar$ gives five numbers,
exactly specified for the $\phi^4$ field, and a fitted kernel reads as physics: the temperature
$\hat\Temp=\hat\bondvar s^2/\latt$, the barrier height and the well position.
\item \emph{Unrestricted and neural kernels} (100 and 1\,516 numbers). A log-linear kernel with a free
two-dimensional spline for $\log\kernel$, which assumes no Gaussian bond, and a mixture density network
(MDN) that maps $a$ to a Gaussian mixture for $a'$. Both are fitted by the same EM; they are the tools for
asking whether a simple family is wide enough (\cref{sec:limits}).
\end{itemize}

\paragraph{Resolution and capacity.}
A flexible kernel can learn narrow features, and the grid must resolve them: with large $\ncomp$ and many
sequences the mixture's narrowest component can be finer than a coarse grid's spacing, which we measure
through a ``resolved'' flag (\cref{sec:compute}). The structure is not free either: $\ncomp=1$ cannot see
fourth-order structure and a large kernel overfits, which is why the tilted and site-energy families have
few parameters.
